Remark that if we choose $t_0\le n^{-\frac{2}{2\beta+d}}$ (as in \cite{oko2023diffusion}),
then the optimal rate may no longer be achieved in \eqref{e18} without further assumptions such as the density lower bound.



\subsection{Analysis of Estimation Error for Diffusion Model}\label{sec_TV_results}
In this section, we analyze the estimation error of the diffusion model in Algorithm~\ref{algorithm1}, based on the score estimation error results from Section~\ref{sec_main_results}. The following theorem concerns the bound of TV distance between the true data distribution $p_0$ and the samples generated by Algorithm~\ref{algorithm1}.
\begin{theorem}\label{main_theorem2}
    Suppose that Assumption~\ref{assumption1} and Assumption~\ref{assumption2} holds. Let $\widehat \bY_{T-t_0}$ be the output of the diffusion model in Algorithm~\ref{algorithm1} at time $T-t_0$. Let $t_0 = n^{-\frac{2}{2\beta+d}}$ and $T = n^{\frac{2\beta}{2\beta+d}}$, then there exists a score estimator $\hat s_t$ such that
    \begin{equation}
    \mathbb{E}\left[\mathrm{TV}\left(\bX_0, \widehat \bY_{T-t_0}\right)\right] \lesssim \mathrm{polylog}(n) \, n^{-\frac{\beta}{2\beta+d}}.
    \end{equation}
\end{theorem}

\begin{remark}
    This sampling error rate coincides with the classical minimax rate in nonparametric density estimation (up to logarithmic factors)
\citep{stone1980optimal,stone1982optimal,tsybakov}, hence must be optimal. The sampling distribution $\widehat \bY_{T-t_0}$ is considered minimax optimal because the sampling error is lower-bounded by the estimation error. To see this, suppose there exists a sampler $\widehat \bZ$ that produces samples with a smaller TV distance $\mathbb{E}\left[\mathrm{TV}\left(\bX_0, \widehat \bZ\right)\right]$ than the minimax rate. We could then generate many independent samples from $\widehat \bZ$ to construct a density estimator for $p_0$ with an error smaller than the minimax rate, leading to a contradiction. Therefore, the sampling error cannot be smaller than the density estimation error.
\end{remark}

\begin{remark}
    In the proof of Theorem~\ref{main_theorem2} (see Section~\ref{subsec_proof_thm2}), we only use the Sobolev class assumption (Assumption~\ref{assumption2}) when controlling the early stopping error $\mathrm{TV}\left(\bX_0, \bX_{t_0}\right)$. Therefore, under only the sub-Gaussian assumption, we can generalize this theorem to other nonparametric classes, provided we can control the early stopping error: $\mathrm{TV}\left(\bX_0, \bX_{t_0}\right)$. Using part 1 of Corollary~\ref{main_corollary} to control the error of score estimation, the sampling error of the diffusion model for a general sub-Gaussian $p_0$ can be bounded as follows:
    \begin{multline*}
        \mathbb{E}\left[\mathrm{TV}\left(\bX_0, \widehat \bY_{T-t_0}\right)\right] \\
